\chapter{Fasting for Tests}

\section*{Definition and Overview}
Fasting before a medical test means eating and drinking nothing except water for a set period beforehand, so that the results reflect your usual health rather than your last meal. Many blood tests are affected by food: glucose and triglycerides rise after eating, and some hormones and other substances vary with meals. A doctor or pathology laboratory will tell you exactly how long to fast, usually 8--12 hours, and whether water, medications, and other drinks are allowed. Getting the fasting right matters, because a test done without proper fasting may give a misleading result and need to be repeated. This chapter explains which tests require fasting, how to prepare, and what to expect.

\section*{Epidemiology}
Fasting blood tests are among the most common laboratory investigations in Australia, with fasting glucose and lipid profiles performed in millions of people each year as part of diabetes and cardiovascular screening. A substantial proportion of fasting samples are taken from people who have not fasted correctly, often because the instructions were unclear, which leads to repeat tests, delayed results, and unnecessary anxiety. Pathology providers report that incorrect fasting is one of the most frequent reasons a test needs to be repeated. Clear communication between clinicians, laboratories, and patients reduces these avoidable repeats.

\section*{Aetiology and Pathophysiology}
The need for fasting arises from how food and drink affect the substances being measured.
\begin{itemize}
    \item \textbf{Glucose:} eating raises blood sugar for hours afterwards, so a non-fasting sample cannot reliably diagnose diabetes or impaired glucose tolerance
    \item \textbf{Lipids:} triglycerides rise markedly after a fatty meal, and cholesterol fractions vary somewhat, so a fasting sample gives a stable baseline for cardiovascular risk
    \item \textbf{Some hormones:} levels of insulin, gastrin, and certain other hormones change with food intake and need a fasting state
    \item \textbf{Other analytes:} iron studies, renal function, and some enzymes can be affected by recent meals in specific circumstances
    \item \textbf{Fluid intake:} water does not interfere and is encouraged to keep veins accessible, while coffee, tea, juice, and alcohol can alter results
    \item \textbf{Medication timing:} some medicines should be taken with food and others must not be stopped, so specific instructions are given for each test
\end{itemize}

\section*{Clinical Features}
Fasting for tests is a procedure rather than an illness, but it has practical features.
\begin{itemize}
    \item A defined fasting window, typically 8--12 hours, during which no food is eaten and only water is drunk
    \item Instructions specific to the test, which may permit water and specify when medications should be taken
    \item The option to schedule the test first thing in the morning so most of the fasting occurs overnight
    \item Possible minor discomforts: hunger, headache, or faintness, especially for people with diabetes or who usually eat breakfast
    \item Clear instructions about whether to continue medicines such as diabetes tablets, insulin, or thyroid replacement
    \item Special considerations for infants and children, for whom fasting is kept as short as safely possible
\end{itemize}

\section*{Differential Diagnosis}
Questions about fasting usually concern how to prepare rather than what is wrong, but considerations include:
\begin{itemize}
    \item \textbf{Which tests require fasting:} not all blood tests do, and fasting unnecessarily can be harmful in some people
    \item \textbf{The exact fasting duration:} this differs between laboratories and tests, so the written instructions should be followed
    \item \textbf{Medication adjustments:} diabetes medicines may need dose changes during fasting, and some medicines must be taken with food
    \item \textbf{Special groups:} pregnant women, children, older adults, and people with diabetes need individualised fasting advice
    \item \textbf{Nutritional safety:} fasting beyond the instructed time, such as overnight to the next day, can cause hypoglycaemia or dehydration
\end{itemize}

\section*{When to Seek Medical Care}
Fasting instructions should be clarified with the doctor or pathology laboratory if anything is unclear, and people with diabetes should discuss how to manage their medicines during the fast.

\textbf{Call 000 (or local emergency services) for:}
\begin{itemize}
    \item Severe hypoglycaemia while fasting: sweating, shaking, confusion, seizures, or loss of consciousness in a person with diabetes
    \item Severe faintness, collapse, or chest pain before or during the blood collection
    \item Any emergency that arises while waiting for test results
\end{itemize}

\section*{Diagnosis and Investigations}
The "investigation" here is the fasting test itself, and preparation determines its validity.
\begin{itemize}
    \item \textbf{Confirm the requirement:} check the pathology request form or ask the laboratory whether fasting applies to your test
    \item \textbf{Time the fast:} count backwards from the appointment so the fast is the correct length, and do not fast for longer than instructed
    \item \textbf{Drink water:} water is allowed and encouraged to keep hydrated and make veins easier to find
    \item \textbf{Plan medications:} ask whether to take regular medicines, especially diabetes and thyroid medicines, and when
    \item \textbf{Reschedule if you break the fast:} if food is eaten by mistake, tell the collector rather than proceeding with an invalid test
    \item \textbf{Special cases:} children and pregnant women may need shorter fasts or modified tests under medical advice
\end{itemize}

\section*{Management}
\begin{itemize}
    \item \textbf{Schedule early:} book the blood test first thing in the morning so most of the fast is spent asleep
    \item \textbf{Keep it simple:} avoid rich meals the night before and stop eating at the instructed time
    \item \textbf{Carry a snack:} bring food to eat immediately after the collection to break the fast promptly
    \item \textbf{Manage diabetes carefully:} follow the specific insulin or tablet plan given by the doctor, and check blood glucose during the fast if advised
    \item \textbf{Communicate honestly:} tell the collector if you did not fast fully so the result can be interpreted correctly
    \item \textbf{Ask about alternatives:} some tests, such as HbA1c for diabetes, do not require fasting and may suit some people better
\end{itemize}

\section*{Complications}
\begin{itemize}
    \item Hypoglycaemia in people with diabetes who fast longer than instructed or adjust medicines incorrectly
    \item Dehydration and faintness from excessive fluid restriction
    \item Headache, irritability, and hunger during the fast
    \item Misleading test results from inadequate fasting, leading to repeat tests, false diagnoses, or inappropriate treatment
    \item Harm from stopping essential medicines without advice
\end{itemize}

\section*{Prognosis and Outcomes}
Fasting for tests is a short, safe procedure for the vast majority of people, and complications are rare when instructions are followed. Correctly fasting ensures that results accurately reflect a person's health, allowing doctors to make reliable decisions about diagnosis and treatment. Where a test result is abnormal or borderline, the clinician will repeat or extend testing as needed. For people with diabetes or other conditions that interact with fasting, a tailored plan agreed with the doctor makes the process safe and straightforward.

\section*{Prevention and Self-Care}
\begin{itemize}
    \item Write down the fasting instructions when the test is ordered and check the pathology request form
    \item Book an early morning appointment to minimise daytime fasting
    \item Drink water freely, but avoid coffee, tea, juice, alcohol, and gum during the fast
    \item Do not smoke on the morning of the test, as smoking can affect some results
    \item Tell the collector about all medicines, supplements, and any mistake in fasting
    \item For children, keep the fast as short as the laboratory allows and follow the doctor's specific advice
    \item Call the laboratory if instructions are missing or unclear
\end{itemize}

\section*{Sources and Further Reading}
The following reputable sources provide further information for healthcare students and patients seeking reliable, current advice about fasting for tests.
\begin{itemize}
    \item Healthdirect Australia. \textit{Blood tests: preparing for your test.} https://www.healthdirect.gov.au/blood-tests
    \item Australian Pathology. \textit{Patient information on fasting.} https://www.pathologyaustralia.com.au/
    \item NHS. \textit{Blood tests: fasting advice.} https://www.nhs.uk/conditions/blood-tests/
    \item Mayo Clinic. \textit{Blood tests: how to prepare.} https://www.mayoclinic.org/
    \item MedlinePlus. \textit{Preparing for a blood test.} https://medlineplus.gov/
    \item MSD Manual. \textit{Common medical tests.} https://www.msdmanuals.com/
\end{itemize}
